\documentclass[11pt]{article}

\usepackage[margin=1in]{geometry}
\usepackage{amsmath,amssymb,amsthm}
\usepackage[ruled,vlined,linesnumbered]{algorithm2e}
\usepackage{mathtools}
\usepackage{xcolor}
\usepackage{booktabs}

\newtheorem{theorem}{Theorem}
\newtheorem{lemma}[theorem]{Lemma}
\newtheorem{corollary}[theorem]{Corollary}
\newtheorem{definition}[theorem]{Definition}
\newtheorem{axiom}[theorem]{Axiom}
\newtheorem{proposition}[theorem]{Proposition}
\theoremstyle{remark}
\newtheorem{remark}[theorem]{Remark}

\DeclareMathOperator{\tid}{tid}
\DeclareMathOperator{\addr}{addr}
\DeclareMathOperator{\kind}{kind}
\DeclareMathOperator{\pc}{pc}
\DeclareMathOperator{\seq}{seq}
\DeclareMathOperator{\loc}{loc}
\DeclareMathOperator{\blk}{block}
\DeclareMathOperator{\scope}{scope}
\DeclareMathOperator{\epoch}{epoch}
\DeclareMathOperator{\Races}{Races}
\DeclareMathOperator{\Writers}{Writers}
\DeclareMathOperator{\Exec}{Exec}
\DeclareMathOperator{\covers}{covers}
\DeclareMathOperator{\val}{val}
\DeclareMathOperator{\rf}{rf}
\DeclareMathOperator{\lab}{lab}

\newcommand{\HB}{\to_{\mathrm{HB}}}
\newcommand{\HBr}{\trianglelefteq_{\mathrm{HB}}}
\newcommand{\CO}{\to_{\mathrm{CO}}}
\newcommand{\PO}{<_{\mathrm{PO}}}
\newcommand{\arr}[2]{\mathrm{arr}^{#2}_{#1}}
\newcommand{\dep}[2]{\mathrm{dep}^{#2}_{#1}}
\newcommand{\ext}[1]{\mathrm{exit}_{#1}}
\newcommand{\tick}[2]{\tau^{#2}_{#1}}
\newcommand{\Fr}{\mathcal{F}}
\newcommand{\Eplus}{E^{+}}

\SetKwSwitch{Switch}{Case}{Other}{switch}{do}{case}{otherwise}{endcase}{endsw}

\title{Happens-Before Data-Race Detection for GPU Kernels:\\
Algorithm and Correctness Proof (v3)}
\date{}

\begin{document}
\maketitle

\noindent\textbf{Changes from v2.}
(1)~Trace validity is no longer an axiom. Well-formedness
(Definition~\ref{def:wf}) is a decidable predicate on traces, stated
declaratively; the implementation's runtime checks are collected in a
separate monitor (Definition~\ref{def:monitor}), and Lemma~\ref{lem:monitor}
proves it rejects exactly the traces that are not well-formed. The single-trace theorems quantify over well-formed traces and assume
nothing about the collector. What the collector must be trusted for is
isolated in Assumptions A1--A3 and used only in
Proposition~\ref{prop:fidelity}, which relates the trace-defined relation to
the execution.
(2)~Barrier records are per-warp arrivals (as emitted); instances are
assembled, not assumed. The clock lemma is restated with tick processing
positions so that it holds in the window between a warp's arrival and its
instance's completion.
(3)~Terminology now follows the verification convention: \emph{sound} = misses
no race, \emph{complete} = reports only races. Theorem~\ref{thm:sound}
(formerly ``completeness'') and Theorem~\ref{thm:complete} (formerly
``soundness'') are renamed accordingly; their content is unchanged.
(4)~\S\ref{sec:cert} is unchanged except for wording and a remark that MM1--MM3
are obligations against the PTX formalisation of Lustig et al., not axioms.
(5)~\S\ref{sec:impl} is updated to the merged state of \texttt{cuVein}.

\paragraph{Terminology.}
Throughout, an analysis is \emph{sound} if it reports a race whenever one
exists (no false negatives) and \emph{complete} if every report is a race (no
false positives).

%=============================================================================
\section{Trace model}\label{sec:model}
%=============================================================================

\begin{definition}[Program, input, execution]
A \emph{program} $P$ is a deterministic kernel: a finite set of threads $t$,
each with a block $\blk(t)$ and a control-flow graph, with no source of
nondeterminism other than memory (randomisation is added in
\S\ref{sec:rand}). An \emph{input} $I$ fixes the initial memory, kernel
parameters and launch configuration, including the block size $N_\beta$ of
each block $\beta$. An \emph{execution} $E$ of $(P,I)$ is one complete run on
the platform; $\Exec(P,I)$ is the set of executions the platform can produce.
\end{definition}

\begin{definition}[Trace records]\label{def:trace}
A trace $T$ is a finite sequence of records; $\seq(e)$ is the position of
record $e$. Records are of three kinds.
\begin{itemize}
\item \emph{Memory records} $e$: $\tid(e)$, $\pc(e)$, $\addr(e)$, a space
      $\in\{\texttt{global},\texttt{shared},\texttt{local}\}$, and
      $\kind(e)\in\{\texttt{read},\texttt{write},\texttt{atomic}\}$. An atomic
      carries $\scope(e)=\mathit{AtomScope}[\pc(e)]\in
      \{\textsc{None},\textsc{Block},\textsc{Grid}\}$,
      $\textsc{None}<\textsc{Block}<\textsc{Grid}$; every atomic PC is in
      $\mathrm{dom}(\mathit{AtomScope})$.
\item \emph{Arrival records} $A$: $\kind(A)\in\{\texttt{barrier},\texttt{syncwarp}\}$,
      a block $\beta(A)$, a warp $\omega(A)$, a lane mask $m(A)$ (the lanes of
      $\omega(A)$ arriving), a barrier index $\mathit{bar}(A)$, and a count
      $n(A)\in\mathbb{N}$ ($0$ meaning ``the whole block''). One record per
      warp per arrival, as emitted by the collector.
\item \emph{Exit records} $\ext{t}$.
\end{itemize}
The \emph{location} of a memory record is
$\loc(e)=(\texttt{shared},\blk(\tid(e)),\addr(e))$,
$(\texttt{local},\tid(e),\addr(e))$ or $(\texttt{global},\addr(e))$.
\end{definition}

\begin{definition}[Expected count, instances]\label{def:instance}
For a \texttt{barrier} arrival $A$ let $\exp(A)=n(A)$ if $n(A)>0$ and
$\exp(A)=N_{\beta(A)}$ otherwise (undefined if $N_{\beta(A)}$ is unknown).
Fix a key $k=(\beta,\mathit{bar})$ and let $A_1,A_2,\ldots$ be the
\texttt{barrier} arrivals with that key in $\seq$ order, with cumulative lane
counts $c_j=\sum_{l\le j}|m(A_l)|$. Provided all $\exp(A_j)$ equal a common
value $x_k$, the $q$-th \emph{segment} of $k$ is
$\{A_j:(q-1)x_k<c_j\le qx_k\}$. A segment whose last element $A_j$ has
$c_j=qx_k$ is \emph{complete} and forms an \emph{instance} $B$ with
participant set $S(B)$ the union (as thread ids) of its lanes and completion
position $\mathrm{cpl}(B)=\seq(A_j)$. A trailing incomplete segment forms no
instance; its arrivals are \emph{unmatched}, with $\mathrm{cpl}=+\infty$.
A \texttt{syncwarp} arrival is an instance by itself, $S(B)$ its masked lanes,
$\mathrm{cpl}(B)$ its own position.
\end{definition}

\begin{definition}[Ghost events, program order]\label{def:ghost}
For each instance $B$ and $t\in S(B)$ introduce an \emph{arrival event}
$\arr{t}{B}$ at the position of the arrival record that covers $t$, and a
\emph{departure event} $\dep{t}{B}$ at position $\mathrm{cpl}(B)^{+}$
(immediately after the completing record; departures of one instance in any
fixed order among themselves). Let $\Eplus$ be the set of memory records,
ghost events and exit records, totally ordered by $\seq$ (extended to ghost
positions). The \emph{program order} $\PO$ of thread $t$ is $\seq$ restricted
to the events of $t$ in $\Eplus$. This is a definition over the trace; it
asserts nothing about the execution.
\end{definition}

\begin{definition}[Atomic successor, scope cover]
For atomics $r,q$ on one address with $\seq(r)<\seq(q)$, $q$ is the
\emph{successor} of $r$ if no atomic on that address lies strictly between
them. With $s=\min(\scope(r),\scope(q))$, $\covers(r,q)$ iff
$s=\textsc{Grid}$, or $s=\textsc{Block}$ and $\blk(\tid(r))=\blk(\tid(q))$.
\end{definition}

\begin{definition}[Happens-before]\label{def:hb}
$\HB$ is the least transitive relation on $\Eplus$ containing
\begin{enumerate}
\item[\textbf{(PO)}] $a\PO b\Rightarrow a\HB b$;
\item[\textbf{(SYNC)}] for every instance $B$ and all $s,t\in S(B)$:
      $\arr{s}{B}\HB\dep{t}{B}$;
\item[\textbf{(ATOM)}] for atomics $r,q$ with $q$ the successor of $r$ and
      $\covers(r,q)$: $r\HB q$.
\end{enumerate}
$a\HBr b$ abbreviates $a\HB b\vee a=b$; $\mathrm{HB}_T$ names the relation.
\end{definition}

\begin{definition}[Conflict, race]\label{def:race}
Memory records $a,b$ \emph{conflict} iff $\loc(a)=\loc(b)$,
$\tid(a)\ne\tid(b)$, and at least one is a \texttt{write} or \texttt{atomic}.
They \emph{race} iff they conflict and neither $a\HB b$ nor $b\HB a$.
$\Races(T)$ is the set of racing pairs.
\end{definition}

\begin{definition}[Well-formed trace]\label{def:wf}
$T$ is \emph{well-formed} iff
\begin{enumerate}
\item[(W0)] the $\seq$ fields of $T$ are strictly increasing along the record
      sequence;
\item[(W1)] for every key $k$, $\exp(A)$ is defined and equal for all
      arrivals $A$ with key $k$; no segment of $k$ overshoots, i.e.\ no $j$
      has $c_{j-1}<qx_k<c_j$ for any $q$; and no thread id occurs twice in one
      segment;
\item[(W2)] for every \texttt{barrier} arrival $A$ and every thread $t$ in
      its lanes, no memory record or arrival record of $t$ lies strictly
      between $A$ and $\mathrm{cpl}(B)$, where $B$ is the instance (or
      unmatched segment, $\mathrm{cpl}=+\infty$) containing $A$.
\end{enumerate}
\end{definition}

W0--W2 are properties of the record sequence alone, decidable in one pass.
(W1) says the instance construction of Definition~\ref{def:instance} closes
exactly; (W2) says a thread does not proceed past a barrier before the
barrier completes. Both are exactly the properties the SYNC rule needs; they
are stated here so that the theorems have a readable hypothesis; a monitor
that decides them is given after Algorithm~\ref{alg:hb}
(Definition~\ref{def:monitor}, Lemma~\ref{lem:monitor}).

\begin{lemma}[HB is a strict partial order below $\seq$]\label{lem:seq}
If $T$ is well-formed, then $a\HB b\Rightarrow\seq(a)<\seq(b)$; in particular
$\HB$ is irreflexive.
\end{lemma}
\begin{proof}
Each generating edge is $\seq$-forward. (PO): by definition of $\PO$ as
$\seq$ restricted to a thread; for the edge $\dep{t}{B}\PO e$ this needs
$\seq(e)>\mathrm{cpl}(B)$, which is (W2). (SYNC):
$\arr{s}{B}$ is at or before $\mathrm{cpl}(B)$, $\dep{t}{B}$ is after it.
(ATOM): by definition of successor. Transitive closure preserves
forwardness.
\end{proof}

\paragraph{Trusted base.}
Nothing above refers to the execution $E$ that produced $T$. The connection is
made once, in Proposition~\ref{prop:fidelity}, under the following assumptions
about the collector. They are the trusted computing base of the tool and are
not checkable from $T$.
\begin{enumerate}
\item[(A1)] \emph{Fidelity.} Every memory access, barrier arrival and thread
      exit of $E$ produces exactly one record, with correct thread, PC,
      address, kind and lane mask, and each thread's records appear in $\seq$
      in its dynamic instruction order.
\item[(A2)] \emph{Coherence order.} For every address, the $\seq$ order of its
      atomics is their coherence order in $E$. (A2 becomes checkable if the
      collector records each RMW's read and written value: the value chain
      determines the coherence order; see \S\ref{sec:impl}.)
\item[(A3)] \emph{Barrier count.} $\exp(B)$ as computed in
      Definition~\ref{def:instance} is the number of threads the hardware
      barrier waits for, and a thread that exits is not among them.
\end{enumerate}

%=============================================================================
\section{Algorithm}\label{sec:alg}
%=============================================================================

\SetKwFunction{Sync}{Sync}
\SetKwFunction{Check}{Check}
\SetKwFunction{Covers}{covers}
\newcommand{\clk}{\mathit{clock}}
\newcommand{\lastR}{\mathit{lastRelease}}
\newcommand{\lastW}{\mathit{lastWrite}}
\newcommand{\rdrs}{\mathit{readers}}
\newcommand{\arrv}{\mathit{arrived}}

\paragraph{State.}
A \emph{clock} is a map from threads to $\mathbb{N}$; a missing entry reads
$0$, and $\sqcup$ is the pointwise maximum. Algorithm~\ref{alg:hb} keeps:
\begin{itemize}
\item $\clk[t]$ --- the vector clock of thread $t$; initially $\clk[t][t]=1$
      and every other entry $0$;
\item $\lastR[x]$ --- for address $x$, the last atomic $r$ on $x$ together
      with the clock $r$ published;
\item $\lastW[\ell]$ --- for location $\ell$, the last write or atomic, as
      (thread, epoch, pc);
\item $\rdrs[\ell]$ --- for location $\ell$, for each thread its latest read
      since $\lastW[\ell]$, as (epoch, pc);
\item $\arrv[k]$ --- the threads arrived so far at the open barrier segment
      of key $k$; initially empty.
\end{itemize}

\begin{algorithm}[H]
\caption{Happens-before oracle with instance assembly}\label{alg:hb}
\KwIn{well-formed trace $T$; $\mathit{AtomScope}$; block sizes $N_\beta$}
\KwOut{set of race reports}
\BlankLine
\SetKwProg{Fn}{function}{:}{}
\Fn{\Sync{$S$}}{
  \tcp{join all participants, then each starts a new epoch}
  $\mathit{joined}\gets$ pointwise maximum of $\clk[t]$ over $t\in S$\;
  \ForEach{$t\in S$}{$\clk[t]\gets\mathit{joined}$;\quad $\clk[t][t]\gets\clk[t][t]+1$\;}
}
\Fn{\Check{$e$}}{
  \tcp{report earlier accesses at $\loc(e)$ not ordered before $e$}
  $t\gets\tid(e)$;\quad $\ell\gets\loc(e)$\;
  \If{$\lastW[\ell]=(w,\mathit{epoch}_w,\mathit{pc}_w)$ \textbf{and} $w\ne t$ \textbf{and} $\mathit{epoch}_w>\clk[t][w]$}{
    report race between $(w,\mathit{pc}_w)$ and $e$\;
  }
  \If{$e$ is a write or an atomic}{
    \ForEach{$(r,(\mathit{epoch}_r,\mathit{pc}_r))\in\rdrs[\ell]$ \textbf{with} $r\ne t$ \textbf{and} $\mathit{epoch}_r>\clk[t][r]$}{
      report race between $(r,\mathit{pc}_r)$ and $e$\;
    }
  }
}
\BlankLine
\ForEach{record $e$ of $T$ in $\seq$ order}{
  \Switch{kind of $e$}{
    \lCase{exit}{do nothing}
    \Case{\texttt{syncwarp} arrival}{
      \Sync{lanes of $e$ as thread ids}\;
    }
    \Case{\texttt{barrier} arrival}{
      $k\gets(\beta(e),\mathit{bar}(e))$;\quad $\arrv[k]\gets\arrv[k]\cup\{\text{lanes of }e\}$\;
      \If{$|\arrv[k]|=\exp(e)$}{
        \Sync{$\arrv[k]$};\quad $\arrv[k]\gets\emptyset$ \tcp*{instance complete}
      }
    }
    \Case{memory access by thread $t$ at location $\ell$}{
      $\mathit{epoch}\gets\clk[t][t]$\;
      \uIf{$e$ is an atomic on address $x$}{
        \tcp{acquire}
        \If{$\lastR[x]=(r,\mathbf{c}_r)$ \textbf{and} \Covers{$r$, $e$}}{
          $\clk[t]\gets\clk[t]\sqcup\mathbf{c}_r$\;
        }
        \Check{$e$}\;
        \tcp{release: publish, then tick}
        $\lastR[x]\gets(e,\clk[t])$\;
        $\lastW[\ell]\gets(t,\mathit{epoch},\pc(e))$;\quad $\rdrs[\ell]\gets\emptyset$\;
        $\clk[t][t]\gets\mathit{epoch}+1$\;
      }
      \uElseIf{$e$ is a write}{
        \Check{$e$};\quad
        $\lastW[\ell]\gets(t,\mathit{epoch},\pc(e))$;\quad $\rdrs[\ell]\gets\emptyset$\;
      }
      \Else{
        \Check{$e$};\quad $\rdrs[\ell][t]\gets(\mathit{epoch},\pc(e))$ \tcp*{read}
      }
    }
  }
}
\Return the reported races\;
\end{algorithm}

\noindent Here \Covers{$r$,$e$} is $\covers(r,e)$ of Definition~\ref{def:hb}.
A race report is the tuple
$(w,\mathit{pc}_w,\tid(e),\pc(e),\loc(e))$. The correspondence with
\texttt{hb\_oracle.py} is: $\clk$ = \texttt{vc}, $\lastR$ = \texttt{released},
$\lastW$ = \texttt{last\_write}, $\rdrs$ = \texttt{last\_reads},
\Sync{} = \texttt{sync\_group}, $\arrv$ = \texttt{pending\_bar}.

\begin{lemma}[Assembly]\label{lem:assembly}
On a well-formed $T$, \Sync{} is invoked exactly once per instance $B$
of Definition~\ref{def:instance}, at position $\mathrm{cpl}(B)$, with
argument $S(B)$, and for no other \texttt{barrier} record.
\end{lemma}
\begin{proof}
For a key $k$, $\arrv[k]$ accumulates the lanes of the open segment
in $\seq$ order and is emptied when $|\arrv[k]|=\exp$. By (W1),
$\exp$ is constant along $k$, no segment overshoots and no thread id repeats
within a segment, so $|\arrv[k]|$ equals the cumulative lane count
$c_j$ modulo $x_k$ and hits $x_k$ exactly at the last arrival of each complete
segment, where $\arrv[k]=S(B)$. A trailing incomplete segment never
triggers the call.
\end{proof}

Algorithm~\ref{alg:hb} is stated for well-formed traces; on other inputs its
behaviour is unspecified. The implementation decides well-formedness online
with the following monitor, which reuses the assembly state and adds two
flags.

\begin{definition}[Well-formedness monitor]\label{def:monitor}
The monitor runs alongside Algorithm~\ref{alg:hb} with additional state
$\mathit{pending}[\omega]\in\{\mathbf{true},\mathbf{false}\}$ per warp
(initially $\mathbf{false}$), $\mathit{expect}[k]$ per key (initially unset),
and $\mathit{prev}$ (initially $-\infty$). At each record $e$ it evaluates the
conditions below, in order, and \emph{rejects} $T$ if any holds; on a
\texttt{barrier} arrival it then sets $\mathit{pending}[\omega(e)]\gets\mathbf{true}$,
$\mathit{expect}[k]\gets\exp(e)$, and, if the instance completes, resets
$\mathit{pending}[\omega]\gets\mathbf{false}$ for every warp with a lane in
$\arrv[k]$; on every record it sets $\mathit{prev}\gets\seq(e)$.
\begin{center}
\begin{tabular}{lll}
\toprule
Record & Reject if & Enforces \\
\midrule
any & $\seq(e)\le\mathit{prev}$ & (W0) \\
\texttt{barrier} arrival & $\exp(e)$ undefined & (W1): $\exp$ defined \\
\texttt{barrier} arrival & $\mathit{expect}[k]$ set and $\ne\exp(e)$ & (W1): $\exp$ constant on $k$ \\
\texttt{barrier} arrival & $\mathit{pending}[\omega(e)]$ & (W1): no repeated lane; (W2) \\
\texttt{barrier} arrival & $|\arrv[k]\cup\text{lanes}(e)|>\exp(e)$ & (W1): no overshoot \\
memory & $\mathit{pending}[\mathrm{warp}(e)]$ & (W2) \\
\bottomrule
\end{tabular}
\end{center}
\end{definition}

\begin{lemma}[Monitor]\label{lem:monitor}
The monitor rejects $T$ iff $T$ is not well-formed.
\end{lemma}
\begin{proof}
The first row rejects iff (W0) fails. Rows two to five reject iff (W1) fails:
$\mathit{expect}[k]$ records $\exp$ of the first arrival on $k$, so row three
fires iff $\exp$ varies along $k$; row five fires iff the running count
crosses a multiple of $x_k$ without landing on it; and a thread id repeated
within a segment is a second arrival of a warp that is still pending, caught
by row four ($\mathit{pending}[\omega]$ is true exactly from a warp's arrival
until its segment completes, by Lemma~\ref{lem:assembly}). Given (W1), rows
four and six reject iff some record of a lane's thread lies strictly between
its arrival and $\mathrm{cpl}(B)$, i.e.\ iff (W2) fails.
\end{proof}

The implementation's monitor differs in two places: it joins per-warp when
$\exp(e)$ is undefined (rejecting only if a second warp arrives in that
state) and it does not check row three. See \S\ref{sec:impl}.

%=============================================================================
\section{Single-trace correctness}\label{sec:single}
%=============================================================================

Throughout this section $T$ is a \emph{well-formed} trace
(Definition~\ref{def:wf}); by Lemma~\ref{lem:monitor} these are exactly the
traces the monitor of Definition~\ref{def:monitor} does not reject. No assumption about the
collector or the platform is used.

\begin{definition}[Ticks, epochs, processing positions]\label{def:epoch}
The \emph{ticks} of thread $t$ are its arrival events and its atomic records,
in $\PO$; $\tick{t}{k}$ is the $k$-th. The \emph{processing position} of a
tick is its own $\seq$ for an atomic and $\mathrm{cpl}(B)$ for
$\arr{t}{B}$. For a memory record $e$ of $t$,
$\epoch(e)=1+\#\{\text{ticks of }t\text{ strictly }\PO e\}$; an atomic $e$
equals $\tick{t}{\epoch(e)}$.
\end{definition}

By (W2) a thread has no memory record between an arrival
and its instance's completion, so the processing positions of $t$'s ticks are
in $\PO$ order and $\epoch(e)$ is also $1$ plus the number of ticks of $t$
whose processing position precedes $\seq(e)$.

\begin{definition}[Frontier]\label{def:frontier}
After the records with $\seq\le i$ have been processed, $\mathrm{last}_t(i)$
is the $\seq$-last event of $t$ in $\Eplus$ at a position $\le i$
(a departure if the last processed record was the completing arrival of an
instance containing $t$; an arrival if $t$ is pending). The \emph{frontier} is
$\Fr_i(t)=\{f\in\Eplus:f\HBr\mathrm{last}_t(i)\}$, empty if $t$ has no event.
For a memory record $b$, $\mathcal{P}(b)=\{f:f\HB b\}$.
\end{definition}

The immediate $\HB$-predecessors of an event are its $\PO$-predecessor; for
$\dep{t}{B}$, the arrivals $\arr{s}{B}$, $s\in S(B)$; for an atomic $q$, its
predecessor $r$ when $\covers(r,q)$. Hence, with $p$ the $\PO$-predecessor of
$b$,
\begin{equation}\label{eq:pred}
\mathcal{P}(b)=\Fr_{\seq(b)-1}(\tid(b))\ \cup\ \{f:f\HBr r\}
\quad(\text{second term only if $b$ acquires from $r$}),
\end{equation}
because $\Fr_{\seq(b)-1}(\tid(b))=\{f:f\HBr p\}$: $p$ is $t$'s last event
before $b$, and by (W2) it is not an arrival of a pending instance.

\begin{lemma}[Clock invariant]\label{lem:vc}
After processing all records with $\seq\le i$:
\begin{enumerate}
\item[(a)] $\clk[t][t]=1+\#\{\text{ticks of }t\text{ with processing position}\le i\}$;
\item[(b)] for $s\ne t$: $\clk[t][s]=\max(\{k:\tick{s}{k}\in\Fr_i(t)\}\cup\{0\})$;
\item[(c)] if $r$ is the $\seq$-last atomic on $x$ with $\seq(r)\le i$, then
      $\lastR[x]=(r,\mathbf{c}_r)$ with
      $\mathbf{c}_r[s]=\max(\{k:\tick{s}{k}\HBr r\}\cup\{0\})$ for every $s$.
\end{enumerate}
\end{lemma}

\begin{proof}
Induction on $i$. Initially all clocks are $\mathbf{0}$ with
$\clk[t][t]=1$, no ticks, empty frontiers; (a)(b) hold and (c) is
vacuous. Let $e$ be the record at position $i$ and assume the invariant at
$i-1$. Two consequences of (a)(b) at $i-1$ are used: for $s\ne t$,
$\clk[t][s]\le\#\{\text{processed ticks of }s\}<\clk[s][s]$, so
$(\bigsqcup_{u\in S}\clk[u])[s]=\clk[s][s]$ whenever $s\in S$
\textup{(\dag)}; and no join raises $\clk[t][t]$.

\emph{Exit.} Nothing changes.

\emph{Read or write $e$ by $t$.} No clock changes and no tick is processed.
$\Fr_i(t)=\Fr_{i-1}(t)\cup\{e\}$ by \eqref{eq:pred}; no tick of any $s\ne t$
is added; other frontiers are unchanged. $c=\clk[t][t]=\epoch(e)$ by
(a) and the remark after Definition~\ref{def:epoch}.

\emph{Non-completing arrival record of warp $\omega$.} No clock changes. For
each covered lane $t$, $\mathrm{last}_t(i)=\arr{t}{B}$ and
$\Fr_i(t)=\Fr_{i-1}(t)\cup\{\arr{t}{B}\}$: the only immediate predecessor of
an arrival is its $\PO$-predecessor. The added event is a tick of $t$, not of
any $s\ne t$, so (b) is unchanged for all threads; (a) is unchanged because
the tick's processing position $\mathrm{cpl}(B)>i$. Other threads' frontiers
are unchanged since the arrival enters no other frontier before
$\mathrm{cpl}(B)$.

\emph{Completing arrival record; instance $B$, $S=S(B)$.} For $t\in S$,
$\mathrm{last}_t(i)=\dep{t}{B}$ and, by the predecessor structure,
\[
\Fr_i(t)=\{\dep{t}{B}\}\ \cup\ \bigcup_{u\in S}\bigl(\{\arr{u}{B}\}\cup\Fr_{a_u}(u)\bigr),
\]
where $a_u$ is the position of $u$'s arrival record; and
$\Fr_{a_u}(u)\setminus\{\arr{u}{B}\}=\Fr_{a_u-1}(u)=\Fr_{i-1}(u)\setminus\{\arr{u}{B}\}$
since $u$ has no event between its arrival and $i$ (W2); that \Sync{}
runs here with exactly $S$ is Lemma~\ref{lem:assembly}. Fix $t\in S$ and $s\ne t$. If $s\notin S$, the
ticks of $s$ in $\Fr_i(t)$ are those in some $\Fr_{i-1}(u)$, $u\in S$, whose
maximal index is $\max_u\clk[u][s]=\mathit{joined}[s]$ by (b). If $s\in S$,
$\arr{s}{B}=\tick{s}{k}$ with $k=\clk[s][s]$ at $i-1$: by (a),
$\clk[s][s]$ counts $1$ plus the ticks processed before $i$, and
$\arr{s}{B}$ is the next tick of $s$ in $\PO$ (no atomic of $s$ lies between
its arrival and $i$); all earlier ticks of $s$ precede it. By \textup{(\dag)},
$\mathit{joined}[s]=k$. So $\clk[t]\gets\mathit{joined}$ establishes (b) for
$t\in S$, and $\clk[t][t]\gets\mathit{joined}[t]+1$ with
$\mathit{joined}[t]=\clk[t][t]$ \textup{(\dag)} establishes (a): exactly one
tick of $t$, $\arr{t}{B}$, has processing position $i$. Threads outside $S$
and (c) are untouched.

\emph{Atomic $q$ by $t$ on $x$.} By (a), $c=\epoch(q)$ and $q=\tick{t}{c}$.
By \eqref{eq:pred}, $\Fr_i(t)=\Fr_{i-1}(t)\cup\{q\}\cup\{f:f\HBr r\}$, the
last term present iff $r$ (the previous atomic on $x$) exists and
$\covers(r,q)$. The algorithm joins $\mathbf{c}_r$ under exactly that
condition, and by (c) at $i-1$, $\mathbf{c}_r[s]$ is the maximal index of a
tick of $s$ in $\{f:f\HBr r\}$; so for $s\ne t$ the joined clock satisfies
(b). The published clock has, for $s\ne t$, the maximal ticks of $s$ in
$\Fr_i(t)\setminus\{q\}=\{f:f\HBr q\}\setminus\{q\}$, and $t$-component $c$,
the index of $q$, which is the $\PO$-last tick of $t$ in $\{f:f\HBr q\}$: this
is (c) for $q$. Then $\clk[t][t]\gets c+1$ restores (a).
\end{proof}

\begin{lemma}[HB test]\label{lem:test}
Let $a$ be a memory record of $s$ and $b$ one of $t\ne s$, $\seq(a)<\seq(b)$.
When \Check{} runs for $b$ (after $b$'s acquire join, if any),
\[
\epoch(a)\le\clk[t][s]\quad\Longleftrightarrow\quad a\HB b .
\]
\end{lemma}
\begin{proof}
By Lemma~\ref{lem:vc}(b)(c) and \eqref{eq:pred}, at that moment
$\clk[t][s]=\max(\{k:\tick{s}{k}\in\mathcal{P}(b)\}\cup\{0\})$.
($\Leftarrow$) On a generating path from $a$ to $b$, the first edge leaving
thread $s$ is (SYNC) from some $\arr{s}{B}$ or (ATOM) from some atomic of
$s$; its source is a tick $\tick{s}{k}$ with $a\HBr\tick{s}{k}$, so
$\epoch(a)\le k$, and $\tick{s}{k}\in\mathcal{P}(b)$ gives
$\clk[t][s]\ge k$. ($\Rightarrow$) Some $\tick{s}{k}\in\mathcal{P}(b)$
with $k\ge\epoch(a)$; every memory record of $s$ with epoch $\le k$ is
$\PO$-before or equal to the $k$-th tick, so $a\HBr\tick{s}{k}\HB b$.
\end{proof}

\begin{theorem}[Completeness: every report is a race]\label{thm:complete}
Let $T$ be well-formed. Every entry $(t_a,p_a,t_b,p_b,\ell)$ that
Algorithm~\ref{alg:hb} adds while processing record $b$ corresponds to a
record $a$ of $t_a$ at $\ell$ with $(a,b)\in\Races(T)$; for a
$\rdrs$ entry, $a$ ranges over the reads of $t_a$ at $\ell$ with the stored
epoch since the last write, all of which race with $b$.
\end{theorem}
\begin{proof}
An entry is added only if $t_a\ne t_b$, both records are at $\ell$, at least
one is a write or atomic, and $\epoch(a)>\clk[t_b][t_a]$ at check
time. By Lemma~\ref{lem:test}, $a\not\HB b$; by Lemma~\ref{lem:seq},
$b\not\HB a$. Reads of $t_a$ at $\ell$ sharing the stored epoch share the
same argument.
\end{proof}

\begin{theorem}[Soundness: every race is reported, location level]\label{thm:sound}
Let $T$ be well-formed. If $(a,b)\in\Races(T)$ with $\seq(a)<\seq(b)$, then
Algorithm~\ref{alg:hb} adds an entry at $\loc(a)$ while processing some
record $b'$ with $\seq(b')\le\seq(b)$.
\end{theorem}
\begin{proof}
Strong induction on $\seq(b)$; $\ell=\loc(a)$, $s=\tid(a)$, $t=\tid(b)$.

\emph{Case A: $a$ is a write or atomic.} When $b$ is processed, \Check{}
compares $b$ with $d=\lastW[\ell]$, present with
$\seq(a)\le\seq(d)<\seq(b)$. If $d=a$ the test fires by
Lemma~\ref{lem:test}. Otherwise $d$ is a write or atomic by $t_d$. If
$t_d\ne t$ and $d\not\HB b$ the test fires for $(d,b)$. Otherwise $d\HB b$
(by (PO) if $t_d=t$; $b\not\HB d$ by Lemma~\ref{lem:seq}). Then
$a\not\HB d$, else $a\HB b$; and $t_d\ne s$, else $a\PO d$. Since
$d\not\HB a$, $(a,d)$ is a race with $\seq(d)<\seq(b)$; apply the induction
hypothesis.

\emph{Case B: $a$ is a read, $b$ a write or atomic.} Let $d$ be the
$\seq$-first write or atomic at $\ell$ with $\seq(a)<\seq(d)\le\seq(b)$
(possibly $d=b$). Until $d$ is processed, $\rdrs[\ell][s]$
holds the epoch of the $\PO$-last read $a_2$ of $s$ at $\ell$ before $d$, with
$a\HBr a_2$. If $a_2\not\HB d$ and $\tid(d)\ne s$, the reader check fires for
$(a_2,d)$ at $d$, $\seq(d)\le\seq(b)$. Otherwise $a\HB d$. If $d=b$ this
contradicts the race. If $d\ne b$: $d\not\HB b$ (else $a\HB b$),
$b\not\HB d$ (Lemma~\ref{lem:seq}), $\tid(d)\ne t$ (else $d\PO b$); so
$(d,b)$ is a race with $d$ a write or atomic, and Case~A applies to it (Case~A
does not use Case~B).
\end{proof}

\begin{remark}[Pair level]
The oracle reports $(a,b)$ itself exactly when $a$ is the $\seq$-last write or
atomic at $\loc(b)$ before $b$, or a $\PO$-last read of its thread at
$\loc(b)$ since that write (the DJIT$^+$/FastTrack guarantee). A verdict
matrix keyed by PC pair must be evaluated at location level, or the oracle run
with full read sets.
\end{remark}

\begin{proposition}[Fidelity]\label{prop:fidelity}
Let $T$ be a well-formed trace of execution $E$ and let $\mathrm{HB}_E$ be the
relation obtained by applying Definition~\ref{def:hb} to the events of $E$
(its dynamic program orders, its barrier instances with their true participant
sets, and its per-address coherence orders). Under A1--A3,
$\mathrm{HB}_T=\mathrm{HB}_E$ on memory records, hence
$\Races(T)=\Races(E)$.
\end{proposition}
\begin{proof}
(PO): by A1, $\seq$ restricted to a thread is its dynamic order.
(SYNC): by A1 each arrival is recorded with its lanes; by A3 an instance
completes exactly when the hardware barrier's participants have all arrived,
so $S(B)$ is the true participant set; departures are placed after
completion, matching the hardware release. (ATOM): successor and $\covers$
depend only on coherence order (A2) and static scopes (A1). The generating
edges coincide, hence so do the closures.
\end{proof}

Proposition~\ref{prop:fidelity} is the only place where the trusted base
enters. Theorems~\ref{thm:complete} and~\ref{thm:sound} are statements about
Algorithm~\ref{alg:hb}, Definition~\ref{def:wf} and $\mathrm{HB}_T$ alone.

\section{From one trace to all executions}\label{sec:cert}
%=============================================================================

\subsection{Why the previous theorem was false}

\begin{remark}[Counterexample]\label{ex:atomic}
Let both atomics be grid-scoped:
\[
t_1:\ x\gets 1;\ \mathrm{atomic}(f)\qquad\qquad
t_2:\ \mathrm{atomic}(f);\ r\gets x .
\]
If $t_1$'s atomic is the coherence predecessor, (ATOM) orders it before
$t_2$'s, so $x\gets1\HB r\gets x$; atomic--atomic pairs are never races, and
the execution is race-free. If $t_2$'s atomic comes first, $r\gets x$ is
unordered with $x\gets 1$: a race. Hence a race-free trace does not certify all
schedules once atomics from different threads meet at an address: the
coherence order is a schedule-dependent input to $\mathrm{HB}_T$, exactly as
the random outcome $R$ is. The certificate below is stated per
\emph{synchronisation profile}.
\end{remark}

\subsection{Memory-model hypotheses}

\begin{definition}[Values, reads-from, profile]\label{def:rf}
Each memory record $e$ of an execution carries $\val_w(e)$ (value written,
for writes and atomics) and $\val_r(e)$ (value read, for reads and atomics),
and a reads-from source $\rf(e)\in\{w:\kind(w)\ne\texttt{read},\ \loc(w)=\loc(e)\}\cup\{\mathit{init}\}$
with $\val_r(e)=\val_w(\rf(e))$ or $I(\loc(e))$.
$\Writers(e)=\{w:\kind(w)\ne\texttt{read},\ \loc(w)=\loc(e),\ w\HB e\}$.
The \emph{profile} $\Pi(E)$ assigns to every address $x$ the sequence of
labels $(\tid,n)$ of the atomics on $x$ in coherence order, where $n$ is the
index of the atomic among that thread's atomics on $x$.
$\Exec(P,I,\Pi)=\{E\in\Exec(P,I):\Pi(E)=\Pi\}$.
\end{definition}

\begin{definition}[Causal order]
$\CO$ is the least transitive relation containing $\HB$ and the \emph{completion
edges} $\ext{s}\CO\dep{t}{B}$ for every barrier instance $B$ in block $\beta$,
every $t\in S(B)$ and every thread $s$ of $\beta$ with $s\notin S(B)$.
\end{definition}

Completion edges express that a block-wide barrier cannot release before every
non-participant of its block has exited; they carry no memory-visibility
guarantee (\texttt{EXIT} has no fence semantics), so they are in $\CO$ and not
in $\HB$.

\begin{axiom}[MM1: causal acyclicity]\label{ax:caus}
For every execution, $\CO\cup\{(\rf(e),e)\}$ is acyclic on $\Eplus$.
\end{axiom}

\begin{axiom}[MM2: coherence]\label{ax:coh}
For every execution and every read or atomic $e$: if $w\HB w'\HB e$ with
$w'\in\Writers(e)$ then $\rf(e)\ne w$; and if $\Writers(e)\ne\emptyset$ then
$\rf(e)\ne\mathit{init}$.
\end{axiom}

\begin{axiom}[MM3: barrier completion]\label{ax:bar}
For every execution and every barrier instance $B$ in block $\beta$, every
thread of $\beta$ either belongs to $S(B)$ or has an exit record; every
\texttt{syncwarp} instance has as participants exactly the lanes of its mask.
Executions terminate.
\end{axiom}

\begin{remark}[Status of MM1--MM3]
MM1--MM3 are not axioms of this document; they are obligations to be discharged
against the published formalisation of the PTX memory model (Lustig,
Sahasrabuddhe, Giroux, \emph{A Formal Analysis of the NVIDIA PTX Memory
Consistency Model}, ASPLOS~2019). The required mapping is: (a) every edge of
Definition~\ref{def:hb} is contained in that model's causality order, which
holds for (SYNC) and holds for (ATOM) only under the fence-gated rule of
\S\ref{sec:impl}, item~4; (b) MM1 and MM2 are instances of that model's
causality and coherence axioms; (c) MM3 is the model's barrier semantics.
Until (a)--(c) are written out, the certificate is a theorem relative to
MM1--MM3, not a fact about hardware.
\end{remark}

\begin{remark}[Model adequacy]\label{rem:adequacy}
The single-trace results of \S\ref{sec:single} are exact with respect to
Definition~\ref{def:hb} and use none of MM1--MM3. The certificate below is only
as good as MM2 is \emph{true on the platform for the edges Definition~\ref{def:hb}
generates}: MM2 asserts that an $\HB$ edge is honoured by the hardware. Under
PTX this holds for (SYNC) edges, but an (ATOM) edge between two
\texttt{.STRONG} atomics is honoured only if the release side is fenced (or
\texttt{.release}) and the acquire side likewise; an unfenced strong atomic is
relaxed. Therefore, in certificate mode, $\mathit{AtomScope}$ must assign
\textsc{None} to an atomic that is not release/acquire-qualified, which is
precisely the fence condition R3 checks with \texttt{release\_scope}. With that
restriction, (ATOM) edges are a subset of PTX synchronises-with and MM1--MM3
are the standard PTX axioms (causality, coherence, barrier semantics). Named
barriers with a partial count have schedule-dependent participant sets and are
outside this section.
\end{remark}

\subsection{First-race prefix and matching}

Fix $P$, $I$, a profile $\Pi$, a race-free $E\in\Exec(P,I,\Pi)$ with well-formed
trace $T$, and any $E'\in\Exec(P,I,\Pi)$ with well-formed trace $T'$. Primed
symbols ($\HB'$, $\CO'$, $\rf'$, $S'$, \dots) refer to $T'$.

\begin{definition}[Labels and counterparts]
Every memory or exit event of a thread $t$ is labelled by $(t,n)$, its index in
$t$'s dynamic order; the $k$-th barrier instance of block $\beta$ is labelled
$(\beta,k)$ (per-warp count for \texttt{syncwarp}), and its ghost events by
$(t,\beta,k,\mathrm{arr}/\mathrm{dep})$. For $e'\in\Eplus(T')$, the
\emph{counterpart} $\varphi(e')$ is the event of $T$ with the same label, if it
exists.
\end{definition}

\begin{lemma}[First race]\label{lem:first}
Suppose $\Races(T')\ne\emptyset$. Let $L'$ be a linear order on $\Eplus(T')$
extending $\CO'\cup\{(\rf'(e),e)\}$ (it exists by MM1). Let $b$ be the
$L'$-least event that is the $L'$-later member of a racing pair, $a$ a racing
partner with $a<_{L'}b$, and $Q=\{e:e<_{L'}b\}$. Then no racing pair of $T'$
has both members in $Q$, and $Q$ is closed under $\CO'$-predecessors.
\end{lemma}
\begin{proof}
A racing pair inside $Q$ would have an $L'$-later member $<_{L'}b$,
contradicting the choice of $b$. Downward closure holds because $L'$ extends
$\CO'$.
\end{proof}

\begin{lemma}[Matching]\label{lem:match}
With $b,Q$ as in Lemma~\ref{lem:first} (or, if $T'$ is race-free,
$Q=\Eplus(T')$ and no $b$), every $e\in Q\cup\{b\}$ satisfies:
\begin{enumerate}
\item[(i)] $\varphi(e)$ exists and agrees with $e$ in kind, pc, location, scope,
      and written value;
\item[(ii)] if $e\in Q$ is a read or atomic, $\val_r(e)=\val_r(\varphi(e))$;
\item[(iii)] (\emph{forward}) if $g\in Q$ and $g\HB' e$ then $\varphi(g)\HB\varphi(e)$;
\item[(iv)] (\emph{backward}) if $g\HB\varphi(e)$ in $T$ then $g=\varphi(g')$ for
      some $g'\in Q$ with $g'\HB' e$.
\end{enumerate}
\end{lemma}

\begin{proof}
Induction along $L'$. Assume the claims for all events $<_{L'}e$; write $t=\tid(e)$.

\emph{Step 1: participants.} Let $B$ be an instance with label $(\beta,k)$ and
$\dep{u}{B}\HBr' e$ for some $u$. We show $S'(B)=S(B_{(\beta,k)})$.
If $s\in S'(B)$ then $\arr{s}{B}<_{L'}e$, so by the induction hypothesis (i)
every event of $s$ up to its $k$-th arrival matches $T$; hence $s$ reaches its
$k$-th barrier in $E$ and $s\in S(B_{(\beta,k)})$. Conversely let
$s\in S(B_{(\beta,k)})\setminus S'(B)$. By MM3 $s$ has an exit in $T'$ and
$\ext{s}\CO'\dep{u}{B}\HBr' e$, so every event of $s$ in $T'$ is
$<_{L'}e$ and matched by (i)(ii): $s$ executes in $T'$ exactly the events it
executes in $T$ up to its exit, with the same read values, hence the same
register state; but in $T$ the same thread continues to a $k$-th arrival. A
deterministic thread cannot both exit and continue from the same state, a
contradiction. The \texttt{syncwarp} case is immediate from MM3.

\emph{Step 2: (i).} The $\PO$-predecessors of $e$ are in $Q$ (as $L'\supseteq\PO$),
and by (i)(ii) for them, $t$ has executed in $T'$ the same events with the same
read values as in $T$. Since $P$ is deterministic, $t$'s next instruction, its
address and any written value are functions of that state; $T$ contains that
next event because $t$ has not exited in $T$ (its exit would be a matched
event). Hence $\varphi(e)$ exists with the claimed agreement. (For ghost events
and exits, (i) is the existence of the counterpart, which follows from the same
argument and Step~1.)

\emph{Step 3: (iii).} A generating edge $g\HB' e$ with $g\in Q$ maps to a
generating edge in $T$: (PO) by labels; (SYNC) by Step~1; (ATOM) because
$\Pi(E)=\Pi(E')$ makes $\varphi(g)$ the coherence predecessor of $\varphi(e)$
with the same scopes and blocks. Chains compose by the induction hypothesis
(iii) on intermediate events, all of which are in $Q$.

\emph{Step 4: (iv).} First let $g$ be an immediate predecessor of $\varphi(e)$
in $T$. If $g$ is the $\PO$-predecessor, its label is that of $e$'s
$\PO$-predecessor $g'\in Q$, and $g'\HB' e$. If $\varphi(e)=\dep{t}{B_{(\beta,k)}}$
and $g=\arr{s}{B_{(\beta,k)}}$ with $s\in S(B_{(\beta,k)})=S'(B)$ (Step~1), then
$g'=\arr{s}{B}\in Q$ and $g'\HB' e$. If $\varphi(e)$ is an atomic acquiring from
$g$, then by equality of profiles $g$'s label is that of $e$'s coherence
predecessor $g'$ in $T'$, which exists in $T'$, satisfies $\covers$, hence
$g'\HB' e$ and $g'\in Q$. For a general $g\HB\varphi(e)$ take a generating path
$g=g_0\to\dots\to g_m=\varphi(e)$; the immediate case gives $g_{m-1}=\varphi(g'_{m-1})$
with $g'_{m-1}\in Q$, $g'_{m-1}\HB' e$; the induction hypothesis (iv) at
$g'_{m-1}$ applied to $g\HB g_{m-1}$ yields $g'\in Q$ with $\varphi(g')=g$ and
$g'\HB' g'_{m-1}\HB' e$.

\emph{Step 5: (ii), for $e\in Q$ a read or atomic.} Let $w'=\rf'(e)$ and
$W=\rf(\varphi(e))$. In $T$, since $E$ is race-free and, by MM1, some linear
order extends $\CO\cup\rf$: if $W\ne\mathit{init}$ then $W$ precedes
$\varphi(e)$ in that order, conflicts with it, and is not racing, so
$W\HB\varphi(e)$; by MM2, $W$ is the $\HB$-maximum of $\Writers(\varphi(e))$
(a chain, as $E$ is race-free), and if $W=\mathit{init}$ then
$\Writers(\varphi(e))=\emptyset$. In $T'$, if $w'\ne\mathit{init}$ then
$w'<_{L'}e$, so $w'\in Q$, $w'$ and $e$ conflict and by Lemma~\ref{lem:first}
do not race, and $e\not\HB' w'$ since $L'$ extends $\HB'$; hence $w'\HB' e$ and,
by MM2, no $u\in\Writers'(e)$ has $w'\HB' u$. If $w'=\mathit{init}$ then
$\Writers'(e)=\emptyset$ by MM2.

Suppose $w'\ne\mathit{init}$. By (iii), $\varphi(w')\in\Writers(\varphi(e))$, so
$W\ne\mathit{init}$ and $\varphi(w')\HBr W$. If $\varphi(w')\ne W$ then
$\varphi(w')\HB W\HB\varphi(e)$; by (iv) at $e$, $W=\varphi(W')$ with $W'\in Q$
and $W'\HB' e$; by (iv) at $W'$, $\varphi(w')=\varphi(u)$ for some $u\in Q$ with
$u\HB' W'$, and $u=w'$ by labels. Thus $w'\HB' W'\HB' e$ with
$W'\in\Writers'(e)$, contradicting MM2 in $T'$. So $W=\varphi(w')$ and
$\val_r(e)=\val_w(w')=\val_w(\varphi(w'))=\val_r(\varphi(e))$ by (i) at $w'$.
Suppose $w'=\mathit{init}$. If $W\ne\mathit{init}$ then by (iv) $W=\varphi(W')$
with $W'\HB' e$, so $\Writers'(e)\ne\emptyset$, contradiction. Hence both read
$I(\loc(e))$.
\end{proof}

\subsection{Certificate}

\begin{theorem}[Single-trace certificate, per profile]\label{thm:cert}
Let $P$ be deterministic and $I$ fixed. If some $E\in\Exec(P,I,\Pi)$ has a
race-free well-formed trace, then every $E'\in\Exec(P,I,\Pi)$ is race-free, and
$T(E')$ consists of the same events as $T(E)$ (same labels, kinds, locations,
values, instances and participants, hence the same $\HB$), up to $\seq$.
\end{theorem}

\begin{proof}
Suppose $\Races(T')\ne\emptyset$ and take $a,b,Q$ from Lemma~\ref{lem:first}.
By Lemma~\ref{lem:match}(i), $\varphi(a)$ and $\varphi(b)$ exist and conflict.
If $\varphi(a)\HB\varphi(b)$, Lemma~\ref{lem:match}(iv) at $b$ gives $a'\in Q$ with
$\varphi(a')=\varphi(a)$ and $a'\HB' b$; labels force $a'=a$, contradicting the
race. If $\varphi(b)\HB\varphi(a)$, (iv) at $a$ gives $b'\in Q$ with
$\varphi(b')=\varphi(b)$, so $b'=b\in Q$, contradicting $b\notin Q$. Hence
$(\varphi(a),\varphi(b))\in\Races(T)$, contradicting race-freedom of $E$.
So $T'$ is race-free; Lemma~\ref{lem:match} with $Q=\Eplus(T')$ gives an
injective label-preserving $\varphi$ that preserves and reflects $\HB$, and it
is surjective because each thread of $E'$ runs to an exit matched in $E$.
\end{proof}

\begin{corollary}[Causal visibility]\label{cor:causal}
If some $E^*\in\Exec(P,I,\Pi)$ has a race then every $E\in\Exec(P,I,\Pi)$ has a
race, and by Theorem~\ref{thm:sound} the oracle reports at least one on any
well-formed trace of any such $E$.
\end{corollary}

\begin{corollary}[Atomic-free kernels]\label{cor:noatom}
If $P$ contains no atomic instruction, $\Pi$ is empty and
$\Exec(P,I,\Pi)=\Exec(P,I)$: a single race-free well-formed trace certifies
race-freedom of $P(I)$ under all schedules, and any racy execution implies
every execution is racy.
\end{corollary}

\begin{remark}
When every address receives atomics from at most one thread, the profile is
forced by program order and the same conclusion is expected; the proof needs an
additional po-loc coherence axiom and is not carried out here.
\end{remark}

\begin{remark}[No iterative completeness]
Fixing a reported race changes $P$; nothing above relates $\Races$ of the
modified program to those of the original. The previous draft's ``iterative
completeness'' is not a theorem.
\end{remark}

%=============================================================================
\subsection{Randomisation}\label{sec:rand}
%=============================================================================

Let a random outcome $R$ fix, per thread, the sequence of values its random
number generator returns. For fixed $(P,I,R)$ the program is deterministic and
everything above holds with $\Exec(P,I,R,\Pi)$ in place of $\Exec(P,I,\Pi)$.
Across different $R$, a branch on a random value changes control flow without
any race, so no race under $R_1$ need have a witness under $R_2$.

\begin{center}
\begin{tabular}{lccc}
\toprule
Property & no atomics, no $R$ & atomics & randomisation \\
\midrule
No false positives on well-formed $T$ (Thm.~\ref{thm:complete}) & \checkmark & \checkmark & \checkmark \\
No false negatives on well-formed $T$ (Thm.~\ref{thm:sound}) & \checkmark & \checkmark & \checkmark \\
One trace certifies $P(I)$ (Thm.~\ref{thm:cert}) & \checkmark & per $\Pi$ & per $(R,\Pi)$ \\
Causal visibility (Cor.~\ref{cor:causal}) & \checkmark & per $\Pi$ & per $(R,\Pi)$ \\
\bottomrule
\end{tabular}
\end{center}

%=============================================================================
%=============================================================================
\section{State of the implementation}\label{sec:impl}
%=============================================================================

Algorithm~\ref{alg:hb} is what \texttt{hb\_oracle.py} and \texttt{HbEngine}
implement on \texttt{cuVein} after the barrier-instance and monitor merges,
with the following differences and open items.

\begin{enumerate}
\item \textbf{Monitor parity.} The implementation joins per-warp when
      $n(A)=0$ and $N_\beta$ is unavailable (rejecting only if a second warp
      arrives in that state), and it takes $\exp$ from the latest arrival
      instead of checking consistency along a key. The monitor of
      Definition~\ref{def:monitor} rejects in both cases, which is what makes
      Lemma~\ref{lem:monitor} an equivalence. Adopt the formal behaviour; on
      every launched kernel $N_\beta$ is known, so the only observable change
      is the added consistency check.
\item \textbf{A2 made checkable.} If the collector records each RMW's read and
      written value, then for each address the coherence order is the unique
      chain in which every RMW reads the value its predecessor wrote, and the
      monitor can verify that $\seq$ agrees with it (M5). This removes A2 from
      the trusted base and turns the emitted profile $\Pi$ into a checked
      fact. Whether the sanitizer callback exposes the operand is to be
      confirmed.
\item \textbf{Profile $\Pi$.} The per-address atomic order and its hash are
      emitted by both artifacts; verdicts are unchanged. Certificates
      (\S\ref{sec:cert}) are per~$\Pi$ and the report should be read that way.
\item \textbf{Fence-gated (ATOM) rule --- open, corpus-affecting.} On the
      target architecture relaxed and release/acquire strong atomics compile to
      the same opcode; only a MEMBAR distinguishes them, and MEMBAR has no
      collector patch point. The implementation therefore treats every strong
      atomic as release/acquire, which fabricates $\HB$ for unfenced atomics
      (a missed-race defect; the discriminating test is tracked as an expected
      failure). The intended fix keeps two separate mechanisms: (i)
      \emph{coherence}: two atomics on one address with covering scope do not
      conflict (a change to Definition~\ref{def:race}, not to $\HB$); (ii)
      \emph{synchronisation}: (ATOM) edges require the releasing PC to be
      fence-qualified before it and the acquiring PC after it, decided
      statically per PC from MEMBAR dominance in the CFG and shipped in the
      sidecar with the coherence scope; effective scope is the minimum of
      atomic and fence scopes. Whether acquire-side qualification is required
      (PTX) or waived (common CUDA practice) is a modelling decision to be
      stated; the strict choice is what maps onto the PTX formalisation. A
      residual window remains: writes between the fence and the atomic are
      published as if released. Until this lands, \S\ref{sec:single} holds for
      the implemented $\HB$ (a legitimate, stronger model) and the certificate
      of \S\ref{sec:cert} is unsound for kernels with unfenced strong atomics.
\item \textbf{Local memory} is keyed by \texttt{(space, addr)} in the
      implementation; Definition~\ref{def:trace} keys it per thread. Align or
      exclude local memory.
\item \textbf{Scale.} The exact oracle's memory is the sum of live clock
      entries; a per-block dense representation with a sparse cross-block
      tail is expected to extend the verified range by an order of magnitude.
      A compact single-epoch engine requires its own lemma (that each
      location's writes are $\HB$-totally ordered in race-free well-formed
      traces under scoped atomics and partial barriers) before it may replace
      Algorithm~\ref{alg:hb} as the reference.
\end{enumerate}

\end{document}
